
\subsubsection{11.11.14}

\begin{enumerate} 
	\item Время начала и окончания собрания:
	17:00 - 20:30
	\item Цели собрания:
	\begin{enumerate}
		\item Завершить работу над механизмом прицепа.
		
		\item Добавить в программу управления роботом управление системой захвата подвижных корзин.
		
		\item Написать программу раздельного управления роботом с двух джойстиков.
		
	\end{enumerate}
	
	\item Проделанная работа:
	\begin{enumerate}
		\item Механизм прицепа был завершен.
		
		\begin{figure}[H]
			\begin{minipage}[h]{0.47\linewidth}
				\center{\includegraphics[scale=0.2]{days/11.11.14/images/01}}  
			\end{minipage}
			\begin{minipage}[h]{0.47\linewidth}
				\center{\includegraphics[scale=0.2]{days/11.11.14/images/02}}
			\end{minipage}
			\caption{Готовый механизм прицепа}
		\end{figure}
		
		\item Программа управления механизмом прицепа не реализована.
		
		\item Сегодня было выбрано окончательное место для NXT-блока. Пока он был временно закреплен7 на скотч, но в будущем мы планируем закрепить его надежнее.
		
		\begin{figure}[H]
			\begin{minipage}[h]{0.2\linewidth}
				\center  
			\end{minipage}
			\begin{minipage}[h]{0.6\linewidth}
				\center{\includegraphics[scale=0.3]{days/11.11.14/images/03}}
				\caption{Место крепления NXT-блока}
			\end{minipage}
		\end{figure}
		
		\item Мы обратили внимание на то, что провод сервопривода, закрепленного на подъемнике, в то время, когда подъемник сложен, касается пола и может помешать движению ковша. Было решено создать специальную катушку, работающую по принципу рулетки и сматывающую провод тогда, когда он не находится в натяжении, либо просто закрепить провод в нескольких местах подъемника таким образом, чтобы он не создавал помех в работе ковша и подъемника.
		
		\item В связи с тем, что нам иногда может потребоваться испытать один из узлов или провернуть с помощью программы приводы, которые трудно провернуть руками, например, для ремонта или замены деталей, было решено создать специальную вспомогательную программу, которая позволила бы управлять отдельными приводами и сервоприводами при помощи кнопок, встроенных в NXT-блок, не прибегая к управлению роботом с джойстика. Такая программа была бы очень удобна в таких случаях, как, например, при внезапной потере связи с роботом по Bluetooth или Samantha, когда необходимо привести механизм подъемника в начальное положение.
		
		\item Программа раздельного управления роботом создана, но не испытана. В новой программе первый оператор отвечает за все, кроме движения, а второй - соответственно за движение робота.
		
		\item На общекомандном обсуждении была разработана идея механизма, служащего для сбивания упора у центральной стойки в автономном режиме: на робота должен быть установлен сервопривод свободного вращения, на котором будут закреплены две цепочки из балок из конструктора LEGO-NXT, скрепленных между собой последовательно таким образом, что каждые две из них соединены между собой только одним штифтом. В сложенном состоянии такая цепочка не будет занимать много места, но с началом вращения, благодаря центробежной силе, распрямится и будет напоминать плеть. После того, как цепь распрямится, для сбития упора центральной стойки роботу будет достаточно проехать от него на расстоянии действия механизма. Это намного проще и удобнее, чем писать программу для поиска упора по ИК-датчику.
		
		\begin{figure}[H]
			\begin{minipage}[h]{0.2\linewidth}
				\center  
			\end{minipage}
			\begin{minipage}[h]{0.6\linewidth}
				\center{\includegraphics[scale=0.2]{days/11.11.14/images/04}}
				\caption{Идея для механизма сбивания упора}
			\end{minipage}
		\end{figure}
		
	\end{enumerate}
	
	\item Итоги собрания: 
	\begin{enumerate}
		\item Механизм прицепа завершен.
		
		\item Программа управления прицепом не реализована.
		
		\item Программа раздельного управления роботом создана.
		
		\item NXT-блок закреплен на роботе.
		
		\item Была разработана концепция механизма сбивания упора.
		
	\end{enumerate}
	
	\item Задачи для последующих собраний:
	\begin{enumerate}
		\item Испытать программу раздельного управления роботом с двух джойстиков.
		
		\item Включить в программу управления роботом программу управления механизмом прицепа.
		
		\item Закрепить провод сервопривода, расположенного на подъемнике так, чтобы он не мешал работе подъемника и ковша.
		
		\item Написать вспомогательную программу для управления узлами робота без использования компъютера и джойстика.
		
		\item Собрать механизм сбивания упора и испытать его в действии.
		
	\end{enumerate}     
\end{enumerate}
\fillpage

